\chapter{Operational continuum and number spectrum (exploratory)}
\label{chap:continuum}
\uses{chap:overview, chap:numbers, chap:algebra}

\section{Position}
\label{sec:continuum_position}

This chapter is \textbf{exploratory} and \emph{not} a published work (no DOI): a
hand-built operational continuum (Path~1, after Brouwer) and an \textbf{operational
number spectrum} $\mathbb{Z} \to \mathbb{Q} \to \mathbb{C} \to \mathbb{R} \to \Omega$,
every node constructed from $\mathbb{Z}$ and kept \textbf{below the
\texttt{Classical.choice} floor}.  Nothing here passes through mathlib's $\mathbb{Q}$
or $\mathbb{R}$, which are entirely Tier-3 (Finding CONT-7: even $(2:\mathbb{Q})+3$
pulls \texttt{Classical.choice}); the construction stays in pure $\mathbb{Z}$
arithmetic, so every object below is \axiomprofile{[propext, Quot.sound]}.

The deliverable is not new mathematics --- these are textbook rationals and Gaussian
rationals --- but the precise \emph{location} of two boundaries, each machine-checked,
and a single thesis they share: the \texttt{Classical.choice} floor is a property of
\emph{construction}, not of the object.

%──────────────────────────────────────────────────────────────────────
\section{The number spectrum}
\label{sec:continuum_spectrum}

\begin{center}
\begin{tabular}{lccc}
\textbf{node} & \textbf{\texttt{DecidableEq}} & \textbf{order} & \textbf{inverse} \\
\hline
$\mathbb{Q}_{\mathrm{op}}$ (\texttt{Qop})    & yes & yes, trichotomy & \textbf{total} \\
$\mathbb{C}_{\mathrm{op}}$ (\texttt{GaussQ}) & yes & none ($\mathbb{C}$ unorderable) & \textbf{total} \\
$\mathbb{R}_{\mathrm{op}}$ (\texttt{Real})   & no (Markov) & apartness only & witnessed (\texttt{Pre.invPos}) \\
\end{tabular}
\end{center}

\noindent Every row is \axiomprofile{[propext, Quot.sound]}.  $\mathbb{Q}_{\mathrm{op}}$
is the \emph{decidable pole}: equality, order and a total reciprocal are all decidable
and choice-free.

\begin{theorem}[Operational $\mathbb{Q}$ is a decidable ordered field-content]
\label{thm:continuum_qop}
\lean{VRCycle.Continuum.Qop.lt_trichotomy, VRCycle.Continuum.Qop.mul_inv_cancel}
\leanok
\texttt{Qop} (cross-multiplication quotient of $\mathbb{Z} \times \mathbb{Z}_{>0}$) is a
\texttt{CommRing} with \texttt{DecidableEq}, full trichotomy
(\texttt{lt\_trichotomy}), and a \emph{total} reciprocal with
$a \ne 0 \to a \cdot a^{-1} = 1$ (\texttt{mul\_inv\_cancel}).  All
\axiomprofile{[propext, Quot.sound]}.
\end{theorem}

\begin{theorem}[Operational $\mathbb{C}$ --- the Gaussian rationals]
\label{thm:continuum_gaussq}
\lean{VRCycle.Continuum.GaussQ.mul_inv_cancel}
\leanok
\uses{thm:continuum_qop}
\texttt{GaussQ} $= \texttt{Qop}[i]$ ($\{re, im : \texttt{Qop}\}$) is a \texttt{CommRing}
with \texttt{DecidableEq} and a total reciprocal $z^{-1} = \bar z/|z|^2$ for $z \ne 0$.
\axiomprofile{[propext, Quot.sound]}.  A \emph{completeness} node: it inherits its base
$\texttt{Qop}$'s operational character and opens no new operational boundary (its lack
of order is classical algebra --- $\mathbb{C}$ is unorderable --- not operationality).
\end{theorem}

%──────────────────────────────────────────────────────────────────────
\section{Boundary I: the Markov line (decidability of zero)}
\label{sec:continuum_markov}

$\mathbb{Q}_{\mathrm{op}}$ and $\mathbb{C}_{\mathrm{op}}$ carry a \emph{total}
reciprocal because zero is \textbf{decidable} there.  Operational $\mathbb{R}$ cannot:
$\lnot(x \approx 0)$ gives no lower-bound modulus on $|x|$ (that step is Markov's
principle, non-constructive), so the operational real reciprocal \texttt{Pre.invPos}
must take an explicit positivity/apartness witness rather than being total.  The
$\texttt{Field}/\texttt{CommRing}$ line across the spectrum is exactly the decidability
of zero.

%──────────────────────────────────────────────────────────────────────
\section{Boundary II: content versus packaging (machine-checked)}
\label{sec:continuum_typeclass}

The field \emph{content} (\texttt{mul\_inv\_cancel}) is choice-free, but registering
\texttt{Qop} as a mathlib \texttt{Field} is blocked below the floor: \texttt{Field}
extends \texttt{DivisionRing extends RatCast}, forcing a map
$\texttt{ratCast} : \mathbb{Q} \to \texttt{Qop}$.  Any such map reads mathlib
$\mathbb{Q}$ (Tier-3), so it pulls \texttt{Classical.choice}.

\begin{finding}[CONT-M5: content free, packaging not]
\label{find:continuum_packaging}
\lean{VRCycle.Continuum.Qop.ofRat}
\leanok
\uses{thm:continuum_qop}
\texttt{Qop.ofRat} (the forced \texttt{ratCast}) is
\axiomprofile{[propext, Classical.choice, Quot.sound]}, while
\texttt{Qop.mul\_inv\_cancel} is \axiomprofile{[propext, Quot.sound]}.  The field's
\emph{doing} (its operations) is operational; the \texttt{Field} \emph{label} (its
packaging) is not.  This is the doing/being asymmetry at the typeclass level.
\end{finding}

\textbf{The boundary is a build invariant.}  The exhibit
\texttt{VRCycle/Continuum/Spectrum.lean} certifies it at build time over the generic
checker \texttt{Meta/DependsOn.lean}:

\begin{center}
\begin{tabular}{lc}
\textbf{declaration} & \texttt{Classical.choice} \\
\hline
\texttt{Qop.mul\_inv\_cancel} (field content)    & free \\
\texttt{GaussQ.mul\_inv\_cancel} (field content) & free \\
\texttt{Qop.lt\_trichotomy} (decidable order)    & free \\
\texttt{Qop.ofRat} (\texttt{ratCast} packaging)  & \textbf{depends} \\
\end{tabular}
\end{center}

%──────────────────────────────────────────────────────────────────────
\section{Two rationals, by necessity}
\label{sec:continuum_two_q}

The repository now contains \emph{two} constructions of the rationals, and this is not
redundancy --- they are the two horns of the floor:
\begin{itemize}
  \item $\mathbb{Q}_{\mathrm{VR}}$ (Chapter~\ref{chap:numbers}, published): built to
  prove an \emph{isomorphism} to mathlib $\mathbb{Q}$; the isomorphism forces the path
  through mathlib and is Tier-3 by necessity.
  \item \texttt{Qop} (this chapter): built to be choice-free below the floor; it makes
  \emph{no} isomorphism claim, which choice-freeness forbids.
\end{itemize}
A single $\mathbb{Q}$ that is \emph{both} choice-free \emph{and} proven isomorphic to
mathlib $\mathbb{Q}$ is impossible: the isomorphism is a \texttt{ratCast}, which pulls
\texttt{Classical.choice} (Finding~\ref{find:continuum_packaging}).  The same object,
two constructions, one above and one below the floor --- the construction-relativity of
the floor, exhibited concretely.

%──────────────────────────────────────────────────────────────────────
\section{Doing, not being (meta)}
\label{sec:continuum_doing}

The same asymmetry appears at the predicate level.  Across the published algebra cycle
the element predicate \texttt{IsOperational} is \texttt{fun \_ => True} in every
instance, and \texttt{Forms.operational\_total} is a \emph{theorem} that operationality
is total: as a property of \emph{being} (the object), operationality does not
discriminate.  What discriminates is the axiom tier of a \emph{construction}, decidable
and universal via \texttt{Meta/DependsOn}.

\begin{theorem}[Operationality: total on being, split on doing]
\label{thm:continuum_doingbeing}
\lean{VR.DoingNotBeing.being_total_int, VR.DoingNotBeing.trans_doing,
      VR.DoingNotBeing.trans_being_via_choice}
\leanok
On one domain $\mathbb{Z}$: \texttt{being\_total\_int} ($\forall a,\ \texttt{IsOperational}\,a$)
is total and choice-free (being does not discriminate), while one $\le$-transitivity
statement has two proofs --- \texttt{trans\_doing} (\texttt{omega})
\axiomprofile{[propext, Quot.sound]} and \texttt{trans\_being\_via\_choice}
(\texttt{le\_trans}) \axiomprofile{[propext, Classical.choice, Quot.sound]} --- so doing
discriminates.  The operational boundary lives on the act, not the object: ``nothing
is, all is doing'' at the meta level.
\end{theorem}

The becoming/Brouwerian core (\texttt{Spread}\,\dots\,\texttt{Model}) realises a
non-enumerable operational continuum via choice sequences
(\texttt{branches\_not\_enumerable}, an unconditional Cantor diagonal,
\axiomprofile{[propext]}), with Brouwerian continuity tracked as a hypothesis or earned
in a finite-information model --- never adopted over classical mathlib, where it is
inconsistent.

%──────────────────────────────────────────────────────────────────────
\section{The three classical pillars, operationally}
\label{sec:continuum_pillars}

The done/becoming cut gives Cantor's diagonal, the Continuum Hypothesis and the Axiom of
Choice a single operational reading.  The machine-checked core is collected here; the
philosophical reading, and the meta-boundaries it cannot cross, are essay-level and marked as
such.  VR does not \emph{solve} these; it \emph{reformulates} them --- operationally they
dissolve or are reinterpreted, formally they remain meaningful.

\begin{theorem}[Operational Cantor]
\label{thm:continuum_cantor}
\lean{VRCycle.Continuum.operational_cantor, VRCycle.Continuum.powerset_diagonal}
\leanok
\uses{thm:continuum_doingbeing}
\texttt{operational\_cantor} \axiomprofile{[propext, Quot.sound]}: the performed (node)
register is countable (\texttt{operational\_register\_countable}, witness \texttt{decodeNode}),
the becoming register is non-enumerable (\texttt{branches\_not\_enumerable}), and the countable
performed cannot exhaust the becoming (\texttt{no\_node\_surjection}).  The $\wp(\mathbb{N})$
bridge \texttt{powerset\_diagonal} \axiomprofile{[propext]}: the full powerset (characteristic
functions $\mathbb{N}\to\mathrm{Bool}$) is non-enumerable, and every enumerable (describable)
family of subsets misses its own diagonal.  So the diagonal \emph{works}, but it cuts
done/becoming --- it does not select an element from a completed uncountable totality (there
is none operationally).  Reinterpreted, not removed.
\end{theorem}

\textbf{CH --- dissolved operationally, retained formally (essay).}  The machine fact under CH
is exactly the $\aleph_0$/non-enumerable gap of \texttt{powerset\_diagonal}.  Beyond it CH adds
no machine content: ``no intermediate cardinal'' is a non-existence claim (classically an
independence result), not provable inside; and the operational layer is not even uniformly
countable (\texttt{Real} is not a countable type --- only \emph{describable} reals would be).
CH is therefore a formal-register label, read philosophically; its honest machine anchor is
\texttt{powerset\_diagonal}.

\begin{theorem}[Operational dependent choice]
\label{thm:continuum_dc}
\lean{VRCycle.Continuum.operational_dependent_choice, VRCycle.Continuum.operational_choice_available}
\leanok
\texttt{operational\_dependent\_choice} \axiomprofile{[propext]}: a history-dependent rule
$f:\mathrm{List\ Bool}\to\mathrm{Bool}$ determines a branch $\alpha\,n = f(\alpha.\mathrm{take}\,n)$
by recursion --- DC is operationally available, choice-free, \emph{because there is a rule}.
Bundled with operational continuity as \texttt{operational\_choice\_available}
\axiomprofile{[propext, Quot.sound]}.  Full AC (selection over an uncountable family with no
rule) has no operational correlate --- a formal-register label; and
$\mathrm{AC}+\mathrm{AD}\to\bot$ is a $T\!\to\!T$ phenomenon that never returns through the
operational floor.  Only DC and WC-N (the \texttt{Continuity} two-register,
\texttt{not\_continuity}\,/\,\texttt{continuity\_of\_nbhd}) are machine objects; full AC and
AC+AD are essay.
\end{theorem}

\textbf{The register asymmetry: $O\to T$ holds, $T\to O$ does not.}  A single machine-checked
asymmetry underlies all three pillars.  In VR-Forms, \texttt{translate\_implies\_realisable}
gives $O\to T$ (an operational fact yields the formal term's realisability); its converse
\emph{fails} --- from $\exists s,\dots$ one cannot recover the specific operational object (no
Skolemisation across the existential).  One passes from operational to formal, but never
\emph{extracts} the operational from the formal: crossing $T\to O$ requires explicit
construction, never derivation.  This is the logical core of register inheritance, and of why
the three pillars dissolve operationally while remaining formally meaningful.

%──────────────────────────────────────────────────────────────────────
\section{Honest scope}
\label{sec:continuum_scope}

\textbf{No new mathematics.}  \texttt{Qop} and \texttt{GaussQ} are textbook
constructions; the value is the choice-free axiom tracking and the two boundaries it
makes precise.  \textbf{$\mathbb{C}$ inherits.}  The Gaussian rationals add a
completeness node, not an operational insight.  \textbf{Choice did not vanish from VR.}
The published works (Numbers, Sets, Audit, Algebra's $\mathbb{Q}$-field) remain Tier-3,
unchanged; this chapter adds a \emph{parallel} choice-free $\mathbb{Q}/\mathbb{C}$
alongside the classical one, not a replacement.  \textbf{The witness is meta-trusted}
--- the dependency checker certifies the relation but is not itself kernel-verified, the
same tier as \texttt{\#print axioms}.

The Lean~4 formalisation lives under \texttt{VRCycle/Continuum/} (\texttt{Rational},
\texttt{GaussianRational}, \texttt{Real}, \texttt{UnitInterval}, the becoming core
\texttt{Spread}\,\dots\,\texttt{Model}, and the consolidated exhibit \texttt{Spectrum}),
with the meta-thesis in \texttt{VRCycle/Meta/DoingNotBeing.lean}.
